\documentclass[11pt]{article}
\usepackage[margin=1in]{geometry}
\usepackage{amsmath,amssymb,amsthm}
\usepackage[colorlinks=true,urlcolor=blue,linkcolor=blue,citecolor=blue]{hyperref}
\newtheorem{theorem}{Theorem}
\newtheorem{lemma}{Lemma}
\newcommand{\Med}{\operatorname{Med}}
\newcommand{\Prob}{\mathbb P}
\newcommand{\E}{\mathbb E}
\title{Strict monotonicity and the limiting median\\of the variance-gamma distribution}
\author{Principia Math in Progress}
\date{12 September 2026}
\begin{document}
\maketitle

\paragraph{Task 1: Resolve MathDB 345410.}
We prove Conjecture 3.5 of Gaunt--Merkle \cite{GM}, including strict
monotonicity, both limiting regimes, and all its median bounds.
The argument compares the distribution function at the decreasing mode
of a law whose gamma shape is one larger.

\begin{theorem}\label{thm:main}
Fix \(r,\theta>0\), and let \(V_{r,\theta,\sigma}\) have the
\(\operatorname{VG}(r,\theta,\sigma,0)\) law. Its median is positive and
strictly decreasing in \(\sigma>0\). If \(q(a)\) denotes the median of
a gamma variable with shape \(a\) and unit rate, then
\begin{align}
 \lim_{\sigma\downarrow0}\Med(V_{r,\theta,\sigma})
   &=2\theta q(r/2),                                       \label{eq:limit0}\\
 \lim_{\sigma\to\infty}\Med(V_{r,\theta,\sigma})
   &=\theta\max\{r-1,0\}.                                  \label{eq:limitinf}
\end{align}
In particular, for every \(\sigma>0\),
\begin{equation}\label{eq:envelopes}
 \theta\max\{r-1,0\}<\Med(V_{r,\theta,\sigma})
                <2\theta q(r/2).
\end{equation}
Consequently,
\begin{align}
 (r-1)\theta
 &<\Med(V_{r,\theta,\sigma})<r\theta e^{-2/(3r)}
   <\left(r-\frac23+\frac{2}{9r}\right)\theta,               \label{eq:bounds}\\
 \Med(V_{r,\theta,\sigma})
 &<(r+2\log2-2)\theta,\qquad r\geq2.                        \label{eq:linear}
\end{align}
The last inequality strengthens the non-strict version in the conjecture.
\end{theorem}

\paragraph{Normalization and a shape-shift identity.}
Let \(a=r/2>0\), \(v=(\sigma/\theta)^2>0\), and \(q_v=\sqrt{1+v}\).
For independent \(X,Y\sim\Gamma(a,1)\), the standard gamma-difference
representation \cite[eq.\ (3.43)]{GM} gives
\begin{equation}\label{eq:representation}
 V_{r,\theta,\sigma}/\theta\ \stackrel d=\ 
 W_{a,v}:=(q_v+1)X-(q_v-1)Y.
\end{equation}
Write \(F_{a,v},f_{a,v}\), and \(M_a(v)\) for its distribution function,
density, and unique median.
Its bilateral Laplace transform, in a neighborhood of zero, is
\[
 L_{a,v}(s)=\E e^{-sW_{a,v}}=(1+2s-vs^2)^{-a}.
\]
For sufficiently small \(s>0\), the transform of its distribution
function is \(L_{a,v}(s)/s\), and hence
\[
 \partial_v\{L_{a,v}(s)/s\}=a\,sL_{a+1,v}(s).
\]
Inversion gives
\begin{equation}\label{eq:pde}
 \partial_vF_{a,v}(x)=a f'_{a+1,v}(x),\qquad x>0.
\end{equation}
To justify this calculation, restrict \(v\) to a compact subset of
\((0,\infty)\) and use the convolution of the two scaled gamma densities.
Parameter derivatives admit integrable bounds with exponential tails;
a common open strip of bilateral Laplace convergence is available.
The shifted density is continuously differentiable, tends to zero at
both infinities, and its derivative has transform \(sL_{a+1,v}(s)\).
Uniqueness of the bilateral transform and continuity give
\eqref{eq:pde} pointwise on \(x>0\).

\paragraph{A Bessel-ratio lemma.}
The auxiliary density \(f_{a+1,v}\) has, for \(x>0\), the form
\[
 f_{a+1,v}(x)=C_{a,v}e^{x/v}x^\nu K_\nu(q_vx/v),
 \qquad \nu=a+\tfrac12,\quad C_{a,v}>0.
\]
The next lemma establishes the parameter behavior of its mode directly.

\begin{lemma}\label{lem:bessel}
For \(\nu>\tfrac12\), put \(d=2\nu-1>0\) and
\[
 R(z)=\frac{K_{\nu-1}(z)}{K_\nu(z)},\qquad
 H(z)=z\left(\frac1{R(z)}-R(z)\right),\qquad z>0.
\]
Then \(R\) increases strictly from \(0\) to \(1\), and \(H\) increases
strictly from \(\max\{d-1,0\}\) to \(d\).
\end{lemma}
\begin{proof}
Use the standard integral
\[
 K_\eta(z)=\int_0^\infty e^{-z\cosh u}\cosh(\eta u)\,du.
\]
The quotient \(R(z)\) is the expectation of
\(g(u)=\cosh((\nu-1)u)/\cosh(\nu u)\) under the probability measure
proportional to \(e^{-z\cosh u}\cosh(\nu u)\,du\).
Since \(|\nu-1|<\nu\), \(g\) is strictly decreasing for \(u>0\):
its logarithmic derivative is
\[
 |\nu-1|\tanh(|\nu-1|u)-\nu\tanh(\nu u)<0.
\]
Thus \(0<R<1\), and differentiation under the integral gives
\[
 R'(z)=-\operatorname{Cov}_z(g(u),\cosh u)>0.
\]
The strict covariance sign follows, for example, by writing the
covariance as half the expectation of
\((g(U)-g(\widetilde U))(\cosh U-\cosh\widetilde U)\)
for independent copies. Standard small- and large-argument expansions
of \(K_\eta\) give the claimed endpoints of \(R\).

The Bessel recurrences give the Riccati equation
\[
 R'=R^2-1+\frac d zR.
\]
Put \(A=(1-R^2)/(1+R^2)\) and \(B=R+1/R\).
Direct differentiation then gives
\begin{equation}\label{eq:Hprime}
 H'=B(H-d+A).
\end{equation}
At any zero of \(H'\), another differentiation gives
\begin{equation}\label{eq:critical}
 H''=BA'=-B\frac{4RR'}{(1+R^2)^2}<0.
\end{equation}
Every critical point would therefore change the derivative from positive
to negative. But the differentiable large-\(z\) expansion is
\[
 R(z)=1-\frac d{2z}+\frac{d(d+2)}{8z^2}+O(z^{-3}),\quad
 H(z)=d-\frac d{2z}+O(z^{-2}),\quad
 H'(z)=\frac d{2z^2}+O(z^{-3})>0
\]
for sufficiently large \(z\). Equations \eqref{eq:Hprime}--\eqref{eq:critical}
exclude a nonpositive value of \(H'\) anywhere: such a value would force
a later change from negative to positive before reaching that final
positive interval. Hence \(H'>0\) everywhere.

Finally, \(H(\infty)=d\). If \(\nu>1\), the small-\(z\) expansion
\(R(z)\sim z/(2\nu-2)\) gives \(H(0+)=2\nu-2=d-1\).
If \(\nu=1\), \(R(z)\sim z\log(1/z)\), giving \(H(0+)=0\).
If \(\tfrac12<\nu<1\), \(R(z)\) is asymptotic to a positive constant
times \(z^{2\nu-1}\), again giving \(H(0+)=0\).
\end{proof}

\paragraph{The mode of the shifted density decreases.}
For \(x>0\), Bessel differentiation gives
\begin{equation}\label{eq:slope}
 \frac{f'_{a+1,v}(x)}{f_{a+1,v}(x)}
 =\frac1v\{1-q_vR(q_vx/v)\},
 \qquad \nu=a+\tfrac12.
\end{equation}
There is a unique positive zero \(D_a(v)\), determined by
\(R(z_v)=1/q_v\) and \(D_a(v)=vz_v/q_v\).
The derivative is positive for \(0<x<D_a(v)\) and negative for
\(x>D_a(v)\). On the negative half-line it is positive, so \(D_a(v)\)
is the unique mode. As \(v\) increases, \(q_v\) increases and \(z_v\)
decreases strictly. At the defining root,
\[
 D_a(v)=z_v(q_v-q_v^{-1})=H(z_v).
\]
Lemma \ref{lem:bessel}, with \(d=2a\), proves
\begin{equation}\label{eq:mode}
 D_a'(v)<0,\qquad
 \lim_{v\downarrow0}D_a(v)=2a,\qquad
 \lim_{v\to\infty}D_a(v)=\ell_a:=\max\{2a-1,0\}.
\end{equation}
In particular \(\ell_a<D_a(v)<2a\).

\paragraph{The median is below this moving mode.}
The representation \eqref{eq:representation} implies
\[
 \frac{W_{a,v}}{q_v}
   =X-Y+\frac{X+Y}{q_v}\longrightarrow X-Y
 \quad\hbox{almost surely as }v\to\infty.
\]
The limiting law is symmetric and has no atom at zero.
Consequently, for every bounded deterministic \(x_v\),
\begin{equation}\label{eq:half}
 F_{a,v}(x_v)\longrightarrow\frac12.
\end{equation}
For example, this follows by squeezing \(x_v/q_v\) between any fixed
positive and negative numbers and then using continuity at zero.

Let \(J(v)=F_{a,v}(D_a(v))\). By \eqref{eq:pde} and the mode equation,
\[
 J'(v)
 =a f'_{a+1,v}(D_a(v))+f_{a,v}(D_a(v))D_a'(v)
 =f_{a,v}(D_a(v))D_a'(v)<0.
\]
The bound \(D_a(v)<2a\) allows \eqref{eq:half} to be applied, so
\(J(v)\to1/2\). Strict decrease then implies \(J(v)>1/2\) for every
finite \(v>0\).

Also \(M_a(v)>0\). Indeed the gamma-difference density satisfies
\(f_{a,v}(x)/f_{a,v}(-x)=e^{2x/v}>1\) for \(x>0\); integration
shows \(F_{a,v}(0)<1/2\), even when the density is unbounded at zero.
We have proved
\begin{equation}\label{eq:medianmode}
 0<M_a(v)<D_a(v).
\end{equation}
Implicit differentiation of the median equation and \eqref{eq:slope}
now give
\begin{equation}\label{eq:Mprime}
 M_a'(v)
 =-a\,\frac{f'_{a+1,v}(M_a(v))}{f_{a,v}(M_a(v))}<0.
\end{equation}
Since \(v=(\sigma/\theta)^2\) is strictly increasing in \(\sigma\),
the desired strict monotonicity follows.

\paragraph{Both limits and all the conjectured bounds.}
When \(a\leq1/2\), \eqref{eq:medianmode} and \eqref{eq:mode} already
give \(M_a(v)\to0=\ell_a\). If \(a>1/2\), then
\(\ell_a=2a-1>0\) is strictly below \(D_a(v)\). Hence
\[
 \partial_v F_{a,v}(\ell_a)=a f'_{a+1,v}(\ell_a)>0.
\]
By \eqref{eq:half}, this increasing function tends to \(1/2\), so
\(F_{a,v}(\ell_a)<1/2\), and
\(\ell_a<M_a(v)<D_a(v)\). Taking \(v\to\infty\) proves
\(M_a(v)\to\ell_a\) in this remaining range.

As \(v\downarrow0\), the same gamma coupling gives
\(W_{a,v}\to2X\) almost surely. The limit has a continuous, strictly
increasing distribution function on its support, so medians converge:
\(M_a(v)\to2q(a)\). Combined with strict monotonicity, this proves
\eqref{eq:limit0}--\eqref{eq:envelopes}.
The standard gamma-median bounds recorded and proved in
\cite[Section 3]{GM} give
\[
 q(a)<a e^{-1/(3a)}\quad(a>0),\qquad
 q(a)\leq a-1+\log2\quad(a\geq1).
\]
Substitution in the strict upper envelope proves
\eqref{eq:bounds}--\eqref{eq:linear}; the last bound in
\eqref{eq:bounds} also uses \(e^{-u}<1-u+u^2/2\) for \(u>0\).
This completes the full conjecture. The equivalent gamma-difference
formulation in \cite[Conjecture 3.3]{GM} is included, with no separate
campaign credit.

\paragraph{Verification and status.}
The companion verifier checks the transform and Bessel differential
identities symbolically and tests the density, mode, and CDF comparisons
numerically. The numerical checks are supplementary to the proof.
The live MathDB claim fields, the current source, its indexed published
text, the 2025 variance-gamma review \cite{FGS}, and later public
literature were checked. No matching full solution claim was located;
the dated audit records the scope and access limitations.

\paragraph{Task 2: Derive rates for the limiting median.}
Find the leading term and a rigorous remainder for
\(M_a(v)-\max\{2a-1,0\}\) as \(v\to\infty\), including the transition
at \(a=1/2\) and any further changes of asymptotic regime.

\paragraph{Task 3: Determine the curvature of the median.}
Classify the signs of \(M_a''(v)\) and of higher variance derivatives,
and determine whether a complete monotonicity statement holds after
subtracting the limiting median.

\paragraph{Task 4: Extend the comparison to unequal gamma shapes.}
Replace \(X,Y\) by independent gamma variables with unequal positive
shapes, choose the positive and negative weights to preserve the mean,
and determine the parameter ranges with a monotone median.

\paragraph{Task 5: Prove the generalized hyperbolic median lower bound.}
Extend the moving-mode argument or its distribution-function identity
to the generalized hyperbolic family and resolve
\cite[Conjecture 3.7]{GM}, including the nonzero mixing-scale parameter.

\begin{thebibliography}{9}
\bibitem{GM}
R.\ E.\ Gaunt and M.\ Merkle,
On bounds for the mode and median of the generalized hyperbolic and
related distributions,
\emph{J.\ Math.\ Anal.\ Appl.} \textbf{493} (2021), 124508.
\href{https://doi.org/10.1016/j.jmaa.2020.124508}{doi:10.1016/j.jmaa.2020.124508}.
\href{https://arxiv.org/abs/2002.01884}{arXiv:2002.01884v4}.
\bibitem{FGS}
A.\ Fischer, R.\ E.\ Gaunt, and A.\ Sarantsev,
The variance-gamma distribution: A review,
\emph{Statistical Science} \textbf{40} (2025), 235--258.
\href{https://doi.org/10.1214/24-STS929}{doi:10.1214/24-STS929}.
\end{thebibliography}
\end{document}
